% Página de créditos (reverso de la portada interior)
\thispagestyle{empty}
\vspace*{\fill}
{\sffamily\small\color{tinta2}\setlength{\parindent}{0pt}\setlength{\parskip}{6pt}
{\titlefont\Large\color{azulnoche} Bioinformática Práctica}\\
{\itshape De la secuencia al sistema: fundamentos, algoritmos y laboratorio computacional}

\medskip
\textbf{Autor:} Juvenal Yosa, PhD · \href{mailto:juvenal.yosa@gmail.com}{juvenal.yosa@gmail.com}\\
\textbf{Copiloto:} Claude (Anthropic), asistente de inteligencia artificial empleado en la
redacción, la programación de figuras y la verificación bibliográfica, bajo la dirección y
revisión del autor.

\textbf{Laboratorios:} cada capítulo tiene notebooks de Google Colab en
\href{https://github.com/juvenalyosa/bioinformatica-practica}{\ttfamily github.com/\allowbreak juvenalyosa/\allowbreak bioinformatica-practica}.

\textbf{Primera edición:} 2026.

\textbf{Licencia:} texto, figuras y código se distribuyen bajo la licencia MIT. Se permite
usar, copiar, modificar y redistribuir este material, con o sin fines comerciales, siempre que se
conserve el aviso de autoría y de licencia.

\textbf{Referencias:} todas las referencias bibliográficas fueron verificadas contra los registros
de Crossref, DataCite u OpenLibrary (DOI o ISBN) y siguen las normas APA (7.ª edición).

\textbf{Tipografía:} compuesto con Lua\LaTeX{} en Libertinus Serif, Source Sans Pro y Fira Mono.
Figuras vectoriales en TikZ y PGFPlots, generadas a partir de cálculos reproducibles.

\medskip
\begin{tikzpicture}
  \foreach \c [count=\i] in {nucA,nucC,nucG,nucT}{\fill[\c] (\i*0.5,0) rectangle ++(0.42,0.12);}
\end{tikzpicture}
\par}
\cleardoublepage
